\documentclass[11pt]{article}
\usepackage{amsmath,amssymb}
\usepackage[hidelinks]{hyperref}
\title{Shared corridor preparation before a delayed density notice}
\date{27 November 2026}
\begin{document}
\maketitle
\paragraph{Audit.}
In the corridor rank-span family of PR~709, let the notice arrive
at a prescribed cut where all robots must be at rest. Both phases
may move every robot while preserving spacing and geometric
link bounds. The first phase is identical in all worlds; the
second uses only the delivered label. The exact required common
preparation depends not just on how many worlds share a label, but
also on where their blocks lie. Nonselected blocks may contract
back toward their hard minimum during the second phase.
The case of observing while still moving is a different problem.

\paragraph{Cut, outputs and budgets.}
There are \(n\ge1\) ordered, contiguous blocks of \(m\ge2\)
adjacent gaps. Initially all gaps are \(g>0\). Every gap must
remain in \([g,R]\), where \(H=m(R-g)\ge\delta>0\). The one hidden
demand at the deadline is a rank-\(m\) span at least
\(mg+\delta\) for the actual block. With a given label function
\(c:\{1,\ldots,n\}\to\{1,\ldots,M\}\), the receivers know the
candidate blocks of that label, not the actual block.
Other than this delivered label there is no distinguishing
information. Temporal evidence costs are not taken to be zero:
the cut is a parameter and this test is conditional on actual
delivery at that cut.

The two rest-to-rest segments have per-robot travel allowances
\(B_1,B_2\), obtained from the motion profile in PR~709, with
unit mobility scales here. Let \(D=2B_2\), and let \(s_j\ge0\)
be the extra span prepared in block \(j\). The total
\(\sum_j s_j\) cannot exceed \(2B_1\), since it equals the net
end-to-end span gain. Conversely every vector with
\(s_j\le H\) and that sum bound can be realized: distribute the
extra span evenly in its block, and center the cumulative
displacements. Initial and terminal hard gaps hold, and the
common normalized profile preserves them between the two cuts.
There is assumed to be enough room at both corridor ends.

\paragraph{Exact residual constraints.}
For a realized label \(h\), let \(z_i\) denote the suffix
displacement of boundary \(i\), shifted by \(B_2\) into
\([0,D]\). The suffix bound for block \(j\) is
\[
 \delta\,{\bf1}_{c(j)=h}-s_j
       \ \le z_j-z_{j-1}\ \le H-s_j .
\]
This follows from the hard minimum and upper span
bounds; every candidate block must meet the additional floor.
Summing lower bounds over an interval \(I=(i,j]\) gives
the exact necessary constraints
\[
 \sum_{k\in I}s_k \ \ge
 w(I):=\left(\delta\max_h |\{k\in I:c(k)=h\}|-D\right)_+ .
\]
For sufficiency, scan the block bounds for a given label with
\(z_0=0\) and \(z_j=\max(0,z_{j-1}+a_j)\),
where \(a_j=\delta{\bf1}_{c(j)=h}-s_j\). Its values are the maximum
of zero and all interval sums of \(a\) ending at \(j\), so the
interval tests bound them by \(D\). If clipping at zero occurs,
the corresponding increment is nonpositive and is still at
least \(a_j\); otherwise it is \(a_j\). In either case it is at
most \(H-s_j\ge0\). Interpolate the boundary displacements
linearly within each block. The middle robots stay within the
same travel allowance and every hard gap in the prepared and
terminal uniformly spaced block is lawful. The common continuous
profile gives an actual suffix motion, not instantaneous restart.

\paragraph{Constructing optimal preparation.}
Write \(p_i=\sum_{j=1}^i s_j\). All interval tests become lower
bounds \(p_j\ge p_i+w((i,j])\). Put \(p_0=0\) and choose
\[
 p_j=\max_{0\le i<j}\{p_i+w((i,j])\}.
\]
This is the least compatible cumulative profile. Since adding
one block increases any interval weight by at most \(\delta\),
\(0\le p_j-p_{j-1}\le\delta\), so the cap \(H\ge\delta\) is
automatically respected. The minimum preparation gain \(p_n\)
is the maximum total weight of a partition into contiguous
intervals, including intervals of zero weight.
The exact feasibility condition for the two phase times and
the given code is \(p_n\le2B_1\).

With \(\delta=D=1\), two notice labels and four blocks,
assigning equal labels to adjacent blocks gives minimum gain 2.
Alternating the labels gives minimum gain 1.
Both codes leave precisely two worlds after a notice, but the
two demands share a middle slack region only in the alternating
case. Treating all previously opened slack as irrevocable would
wrongly forbid this redistribution.

\paragraph{Optimal label design.}
Among codes with at most \(M\ge1\) labels, a cyclic assignment
\(c(j)=1+((j-1)\bmod M)\) minimizes every interval weight at once.
Any labeling of an interval of length \(\ell\) has a label
with at least \(\lceil\ell/M\rceil\) occurrences; cyclic ones
attain it simultaneously. Hence the exact optimized minimum
preparation gain is
\[
 \begin{cases}
  (\delta\lceil n/M\rceil-D)_+,&D\ge\delta,\\[2mm]
  \displaystyle\max_{1\le p\le n}
       \left[p(\delta-D)+\delta\left\lfloor(n-p)/M\right\rfloor\right],
       & D<\delta .
 \end{cases}
\]
For the first branch the cyclic interval weight cannot decrease
in total when adjacent parts are merged: two positive weights lose
at most one \(\delta\) in their combined ceiling count but pay
the suffix allowance \(D\) only once. If one weight is zero,
monotonicity is enough. For the second branch every nonempty part
is positive; for \(p\) parts the sum of ceiling counts is at most
\(p+\lfloor(n-p)/M\rfloor\), attained by giving \(p-1\) parts
one block each and putting the rest in a single part.

\paragraph{Boundary and verification.}
The formula is for a mandatory rest cut, not a restriction that
may silently be imposed on all possible controllers. With three
blocks, no informative notice, \(\delta=1\), unit scales,
\(a=2\), large speed bound and \(T=2\), uninterrupted
rest-to-rest motion has budget 2, sufficient for all three
blocks. Splitting at time 1 into two rest-to-rest phases gives
budget \(1/2\) in each phase, insufficient. Momentum across the
cut can therefore change the yes/no result. This negative control
does not solve the general moving-notice problem.

An independent integer enumeration of all prefix gains and
boundary suffix motions checks 4,300 scaled queries; 150 tests
compare best codes, cyclic assignment and the closed form.
Four Lean lemmas cover the tight displacement, capacity-clipping,
centered witness, and interval-merge steps. A direct matched
difference/interval solver can express and attain the same
solutions. These are model guarantees, not a source of the
delayed message or a deployment result.
\end{document}
